\documentclass[11pt, a4paper]{../problems_template}

\begin{document}

\begin{titlepage}
  \begin{center}
	\Huge{problems.ru}
  \end{center}
\end{titlepage}   

\begin{exercise}
В банде 50 бандитов. Все вместе они ни в одной разборке ни разу не участвовали, а каждые двое встречались на разборках ровно по разу. Докажите, что один из бандитов был не менее, чем на восьми разборках.
\end{exercise}

\begin{Solution}
Предположим противное. Выберем произвольного бандита, полагая, что он был менее, чем на 8 разборках. С каждым из 49 оставшихся он должен встретиться не более раза, значит, на какой то разборке он должен был встретиться не менее чем с 7 бандитами. Всего на этой разборке было не менее 8 бандитов. Выберем бандита, который в этой разборке не принимал участие. Он искомый, так как с каждым из бандитов, бывших на этой разборке он должен был встретиться по разу. 
\end{Solution}

\begin{exercise}
За круглым столом совещались $2n$ депутатов. После перерыва эти же $2n$ депутатов расселись вокруг стола, но уже в другом порядке. Доказать, что найдутся два депутата, между которыми как до, так и после перерыва сидело одинаковое число человек.
\end{exercise}

\begin{exercise}
Дракон запер в пещере шестерых гномов и сказал: "У меня есть семь колпаков семи цветов радуги. Завтра утром я завяжу вам глаза и надену на каждого по колпаку, а один колпак спрячу. Затем сниму повязки, и вы сможете увидеть колпаки на головах у других, но общаться я вам уже не позволю. После этого каждый втайне от других скажет мне цвет спрятанного колпака. Если угадают хотя бы трое, всех отпущу. Если меньше - съем на обед". Как гномам заранее договориться действовать, чтобы спастись?
\end{exercise}

\begin{Solution}
\end{Solution}

\begin{exercise}
Докажите, что среди любых шести человек есть либо трое попарно знакомых, либо трое попарно незнакомых.
\end{exercise}

\begin{Solution}
Рассмотрим раскраску полного графа на $6$ вершинах в черный и белый цвета. Выберем произвольную вершину $A$. Ее степень равна 5, а значит среди её ребер найдутся три одного цвета. Пусть этот цвет черный. Обозначим эти ребра $AB_{1}$, $AB_{2}$, $AB_{3}$. Тогда либо одно из ребер $B_{1}B_{2}$, $B_{2}B_{3}$, $B_{1}B_{3}$ будет черным, и у нас будет полный подграф черного цвета, либо все три ребра белые, и тогда будет полный подграф белого цвета.
\end{Solution}

\begin{exercise}
Каждое из рёбер полного графа с $17$ вершинами покрашено в один из трёх цветов. Докажите, что есть три вершины, все рёбра между которыми - одного цвета.
\end{exercise}

\begin{Solution}
Рассмотрим раскраску полного графа на $17$ вершинах в красный, зеленый и синий цвета. Выберем произвольную вершину $A$. Среди её рёбер найдётся не менее шести рёбер одного цвета - пусть это будет красный. Тогда либо среди рёбер, соединяющих вершины, смежные с $A$ по красным рёбрам, найдётся красное ребро (и мы получим полный красный подграф), либо образуется полный подграф с рёбрами двух цветов. А значит, в нём найдётся полный подграф зелёного или синего цвета.
\end{Solution}

\begin{exercise}
Докажите, что для плоского связного графа справедливо неравенство  $E \leqslant  3V - 6$.
\end{exercise}

\begin{Solution}
Любое ребро делит две грани, а любая грань состоит не менее чем из трех ребер. То есть $2E \geqslant 3G$. Тогда 
$$
V - E + G = 2 \implies 3V - 3E + 3G = 6 \implies 3V - 3E + 2E \geqslant 6 \implies E \leqslant 3V - 6
$$  
\end{Solution}

\begin{exercise}
Каждое ребро полного графа с $11$ вершинами покрашено в один из двух цветов: красный или синий. Докажите, что либо красный, либо синий граф не является плоским.
\end{exercise}

\begin{Solution}
Всего в графе $\binom{11}{2} = 55$ ребер. Значит ребер какого-то цвета не менее 28. Тогда для подграфа с ребрами этого цвета $3V - 6 = 27 \leqslant 28 \leqslant E$. Значит этот подграф не планарен.
\end{Solution}

\begin{exercise} 
Каждый из учеников класса занимается не более чем в двух кружках, причём для любой пары учеников существует кружок, в котором они занимаются вместе. Докажите, что найдётся кружок, в котором занимается не менее $\frac{2}{3}$ всего класса.
\end{exercise}

\begin{Solution}

\end{Solution}

\begin{exercise}
В графе 20 вершин, степень каждой не меньше 10. Доказать, что в нём есть гамильтонов путь.
\end{exercise}

\begin{Solution}
Допустим противное. Тогда рассмотрим самый длинный из существующих путей, обозначив его $L$. Пусть в $L$ входит $n$ вершин. Вершины которые не входят в $L$ обозначим $M$. Вершины пути $L$ будем обозначать $L_{i}$ где $1 \leqslant i \leqslant n$. Так как степень $L_{n}$ не меньше 10 и это конечная вершина, то все ребра из нее должны идти в вершины $L$. Если из вершины $L_{i}$ есть ребро в вершину из $M$ и существует ребро из $L_{n}$ в $L_{i - 1}$ то мы можем сконструировать путь $L_{1} \dots L_{i-1}L_{n}\dots L_{i} M$. Длина этого пути будет больше $n$. Отсюда следует, что в $M$ могут идти только ребра тех вершин $L_{i}$, которые не соединены с $L_{n}$. Таких ребер не более $n - 10$. В компоненте $M$ не более $20 - n$ вершин, то есть степень вершины в $M$ не больше $19 - n$. Из $L$ в $M$ ведут не более $n - 10$ ребер. Значит максимальная степень вершины в $M$ равна $(19 - n) + (n - 10) = 9$. Противоречие.
\end{Solution}

\begin{exercise}
Каждое из рёбер полного графа с 9 вершинами покрашено в синий или красный цвет. Докажите, что либо есть четыре вершины, все рёбра между которыми - синие, либо есть три вершины, все рёбра между которыми - красные.
\end{exercise}

\begin{Solution}
  
\end{Solution}

\begin{exercise}
Каждое из рёбер полного графа с 18 вершинами покрашено в один из двух цветов. Докажите, что есть четыре вершины, все рёбра между которыми - одного цвета.
\end{exercise}

\begin{Solution}
Рассмотрим произвольную вершину. По принципу Дирихле, от нее отходят не менее 9 ребер одного цвета. Пусть этот цвет красный. В графе на 9 вершинах, к которым идут эти ребра, найдется три вершины, все ребра между которыми - красные или четыре вершины все ребра между которыми синие. В каждом случае мы получим клику на 4 вершинах, все ребра между которыми одного цвета.  
\end{Solution}


\begin{exercise}
Дно прямоугольной коробки вымощено плитками $1 \times 4$ и $2 \times 2$. Плитки высыпали из коробки и одна плитка $2 \times 2$ потерялась. Ее заменили на плитку $1 \times 4$. Докажите, что теперь дно коробки вымостить не удастся.
\end{exercise}

\begin{Solution}
Рассмотрим раскраску вида 
$$
\begin{matrix}
1 & 2 & 1 & 2 & \dots \\
3 & 4 & 3 & 4 & \dots \\
1 & 2 & 1 & 2 & \dots \\
3 & 4 & 3 & 4 & \dots \\
\vdots & \vdots & \vdots & \vdots & \ddots
\end{matrix}
$$
Любая плитка $2 \times 2$ содержит 1 клетку цвета 1. Значит четность количества плиток этого вида должна быть равна четности количества клеток цвета 1. Так как плитка вида $1 \times 4$ не содержит вовсе или содержит 2 клетки цвета 1, то замощение при замене плитки выполнить не получится.
\end{Solution}

\end{document}
